\documentclass[12pt, a4paper]{article}

\usepackage[french]{babel}
\usepackage{geometry}
\usepackage{hyperref}
\usepackage{listings}
\usepackage{xcolor}
\usepackage{booktabs}
\usepackage{tabularx}
\usepackage{array}
\usepackage{tikz}
\usetikzlibrary{shapes.geometric, arrows.meta, positioning}
\usepackage{amsmath}
\usepackage{enumitem}
\usepackage{fancyhdr}
\usepackage{titlesec}
\usepackage{mdframed}
\newcommand{\up}[1]{\textsuperscript{#1}}

\geometry{margin=2.5cm}

% En-têtes
\pagestyle{fancy}
\fancyhf{}
\rhead{Agent Actuariat — Architecture LLM}
\lhead{\leftmark}
\cfoot{\thepage}

% Couleurs
\definecolor{codegreen}{rgb}{0,0.6,0}
\definecolor{codegray}{rgb}{0.5,0.5,0.5}
\definecolor{codepurple}{rgb}{0.58,0,0.82}
\definecolor{backcolour}{rgb}{0.97,0.97,0.97}
\definecolor{alertblue}{rgb}{0.0,0.3,0.7}
\definecolor{alertyellow}{rgb}{0.9,0.7,0.0}

% Style de code
\lstdefinestyle{pythonstyle}{
    backgroundcolor=\color{backcolour},
    commentstyle=\color{codegreen},
    keywordstyle=\color{codepurple},
    numberstyle=\tiny\color{codegray},
    stringstyle=\color{codegreen},
    basicstyle=\ttfamily\footnotesize,
    breakatwhitespace=false,
    breaklines=true,
    captionpos=b,
    keepspaces=true,
    numbers=left,
    numbersep=5pt,
    showspaces=false,
    showstringspaces=false,
    showtabs=false,
    tabsize=2,
    frame=single,
    rulecolor=\color{codegray}
}

\lstset{style=pythonstyle}

% Boîte d'alerte
\newmdenv[
  backgroundcolor=backcolour,
  linecolor=alertblue,
  linewidth=2pt,
  leftmargin=0,
  rightmargin=0,
  innerleftmargin=10pt,
  innerrightmargin=10pt,
  innertopmargin=8pt,
  innerbottommargin=8pt,
]{infobox}

\newmdenv[
  backgroundcolor=yellow!10,
  linecolor=alertyellow,
  linewidth=2pt,
  leftmargin=0,
  rightmargin=0,
  innerleftmargin=10pt,
  innerrightmargin=10pt,
  innertopmargin=8pt,
  innerbottommargin=8pt,
]{warningbox}

% Titre
\title{
  \Large\textbf{Agent Actuariat}\\[0.5em]
  \large Réflexion architecturale et refactorisation\\
  vers une gestion dynamique du contexte LLM\\[0.5em]
  \normalsize Migration vers LangGraph et catalogue à 3 niveaux
}
\author{Marc Juillard}
\date{31 mars 2026}

\begin{document}

\maketitle
\tableofcontents
\newpage

% ─────────────────────────────────────────────────────────────────────────────
\section{Contexte et problème initial}
% ─────────────────────────────────────────────────────────────────────────────

\subsection{Le problème constaté}

Lors de l'exécution du pipeline de construction de table de mortalité sur le
portefeuille \texttt{données\_actif\_decede.txt} (530~345 lignes, 7 colonnes),
le pipeline crashait systématiquement à partir de l'étape~19 avec l'erreur
suivante :

\begin{warningbox}
\texttt{Error code: 429 --- Limit 30,000 TPM, Requested 30,440. Request too
large for gpt-4o.}
\end{warningbox}

Ce crash intervenait à l'étape du benchmarking, juste avant la génération du
rapport PDF, rendant le pipeline incapable de produire son livrable final.

\subsection{Anatomie d'un appel LLM dans le pipeline}

Un pipeline de construction de table de mortalité comprend une dizaine d'étapes
(exposition, taux bruts, diagnostics, lissage, validation, benchmarking,
rapport PDF). À chaque étape, l'agent envoie à l'API OpenAI un message
contenant \emph{la totalité du contexte accumulé} depuis le début de la
session. En effet, les modèles de type GPT ne maintiennent aucune mémoire
inter-appels : chaque requête est autonome et doit embarquer l'intégralité
de l'historique conversationnel.

\begin{figure}[h]
\centering
\begin{tikzpicture}[
  box/.style={rectangle, draw, rounded corners, minimum width=8cm, minimum height=0.7cm, fill=backcolour},
  arrow/.style={->, thick, >=Latex}
]
\node[box] (t1) at (0, 6) {Tour 1 : prompt système (17k) + message user};
\node[box] (t2) at (0, 5) {Tour 2 : prompt (17k) + user + résultat exposure};
\node[box] (t3) at (0, 4) {Tour 3 : prompt (17k) + user + exposure + crude\_rates};
\node      (dots) at (0, 3) {\vdots};
\node[box, fill=red!15] (t10) at (0, 2) {Tour 10 : \textbf{30~440 tokens} $\rightarrow$ \textbf{429}};
\draw[arrow] (t1) -- (t2);
\draw[arrow] (t2) -- (t3);
\draw[arrow] (t3) -- (dots);
\draw[arrow] (dots) -- (t10);
\end{tikzpicture}
\caption{Croissance linéaire du contexte LLM à chaque tour du pipeline}
\end{figure}

La taille du contexte croît à chaque tour d'environ 500 à 800 tokens
(résultat du tool call précédent). Avec un prompt système de 17~000 tokens,
la limite de 30~000 TPM est atteinte dès le dixième appel.

\subsection{La cause racine}

L'architecture précédente injectait un catalogue enrichi dans le prompt
système. Ce catalogue, généré à partir des contrats de 19 tools actuariels
(11 sections documentaires chacun), pesait \textbf{17~000 tokens}. Il était
retransmis \emph{intégralement à chaque appel LLM}, soit environ 10 fois pour
un pipeline complet.

\begin{table}[h]
\centering
\begin{tabular}{lrr}
\toprule
Composant & Avant refactorisation & Après refactorisation\\
\midrule
Prompt système & 17~000 tokens & 4~553 -- 11~595 tokens\\
Coût total (12 tours) & $\sim$230~000 tokens & $\sim$90~000 tokens\\
Coût estimé par run & $\sim$0,60 USD & $\sim$0,25 USD\\
Crash au tour & 10 & jamais\\
\bottomrule
\end{tabular}
\caption{Impact de la refactorisation sur la consommation de tokens}
\end{table}

% ─────────────────────────────────────────────────────────────────────────────
\section{Réflexion sur la nécessité du catalogue}
% ─────────────────────────────────────────────────────────────────────────────

\subsection{Tous les champs sont-ils utiles à chaque appel ?}

La première question posée était la suivante : \emph{l'agent a-t-il besoin de
tout le catalogue à chaque appel LLM ?}

Une analyse champ par champ du catalogue YAML a conduit à distinguer trois
catégories d'information :

\begin{table}[h]
\centering
\begin{tabularx}{\textwidth}{lXl}
\toprule
Champ & Utilité & Verdict\\
\midrule
\texttt{short\_description} & Décider quel tool appeler & Indispensable\\
\texttt{depends\_on} & Construire l'ordre d'exécution & Indispensable\\
\texttt{when\_to\_use / not} & Éviter les appels hors contexte & Indispensable\\
\texttt{prerequisites} & Vérifier ce qui doit exister avant & Indispensable\\
\texttt{quality\_gates.blocking} & Règles comportementales critiques & Indispensable\\
\texttt{reasoning\_hint} & Guidance directe de l'agent & Indispensable\\
\texttt{params} & Paramètres à passer à la fonction & Indispensable\\
\midrule
\texttt{outputs.return\_payload} & L'agent voit les vraies clés en résultat & Bruit\\
\texttt{description} & Version longue de \texttt{short\_description} & Redondant\\
\texttt{display\_name} & Usage UI uniquement & Inutile pour l'agent\\
\texttt{domain / capability\_group} & Métadonnées organisationnelles & Inutile pour l'agent\\
\texttt{required\_by} & Inverse de \texttt{depends\_on}, superflu & Inutile pour l'agent\\
\texttt{exemplar\_query} & Redondant avec \texttt{reasoning\_hint} & Supprimable\\
\texttt{client\_visible} & Usage UI uniquement & Inutile pour l'agent\\
\bottomrule
\end{tabularx}
\caption{Classification des champs du catalogue selon leur utilité réelle pour l'agent}
\end{table}

\subsection{La question clé : quand chaque information est-elle nécessaire ?}

L'analyse plus fine a révélé que la valeur du catalogue varie non seulement
selon les champs, mais aussi selon la \emph{phase cognitive} de l'agent :

\begin{infobox}
\textbf{Phase de qualification} : l'agent doit répondre à ``est-ce faisable ? qu'est-ce qu'il me manque ?''. Il n'a besoin que de la description courte
et des prérequis.

\textbf{Phase de planification} : l'agent construit la séquence complète avec
paramètres. C'est le seul moment où la totalité des informations est exploitée.

\textbf{Phase d'exécution} : l'agent avance pas à pas. Il connaît déjà son plan
et n'a besoin que des quality gates pour détecter une anomalie.
\end{infobox}

% ─────────────────────────────────────────────────────────────────────────────
\section{Architecture à 3 niveaux de catalogue}
% ─────────────────────────────────────────────────────────────────────────────

\subsection{Les 3 niveaux}

La solution retenue consiste à faire varier le contenu du prompt système selon
la phase cognitive active, via trois niveaux de catalogue :

\begin{table}[h]
\centering
\begin{tabularx}{\textwidth}{llXl}
\toprule
Niveau & Tokens & Champs injectés & Déclencheur\\
\midrule
MIDDLE & $\sim$2~500 & \texttt{short\_description}, \texttt{when\_not\_to\_use},
  \texttt{prerequisites} & 1\textsuperscript{er} message user, plan non établi\\
FULL & $\sim$8~000 & Tous les champs utiles & Message user + plan existant (replanification)\\
LIGHT & $\sim$1~300 & \texttt{short\_description}, \texttt{depends\_on},
  \texttt{quality\_gates.blocking} & Dernier message = résultat de tool\\
\bottomrule
\end{tabularx}
\caption{Les 3 niveaux de catalogue et leurs déclencheurs}
\end{table}

\subsection{Logique de transition}

Le niveau actif est déterminé à chaque tour par une règle simple sur le type
du dernier message dans la conversation :

\begin{lstlisting}[language=Python, caption={Logique de sélection du niveau de catalogue}]
def _get_catalogue_level(state: AgentState) -> str:
    last_msg = state["messages"][-1]

    # Tour d'exécution : dernier message = résultat de tool
    if isinstance(last_msg, ToolMessage):
        return "light"

    # Premier contact : l'agent ne sait pas encore ce qu'il peut faire
    if not state["plan_established"]:
        return "middle"

    # Message utilisateur avec plan existant = replanification
    return "full"
\end{lstlisting}

La variable \texttt{plan\_established} passe à \texttt{True} dès le premier
tool call exécuté.

\subsection{Cas d'un retour à FULL en cours d'exécution}

Un point important a été identifié : l'exécution n'est pas toujours linéaire.
Trois situations justifient un retour au niveau FULL en cours de pipeline :

\begin{enumerate}
  \item Un tool retourne une erreur structurelle (colonne manquante, données
    corrompues) qui invalide l'ensemble du plan.
  \item Un résultat remet en cause les hypothèses initiales (ex.~: 80\% des
    âges non crédibles, le pipeline standard ne s'applique pas).
  \item L'utilisateur envoie un nouveau message en cours d'exécution,
    réorientant volontairement l'analyse.
\end{enumerate}

En revanche, les ajustements mineurs (doubler \texttt{lambda\_wh} si
\texttt{n\_non\_monotone} > 0) \emph{ne} justifient pas un retour à FULL :
ils sont gérés par les quality gates dans le catalogue LIGHT et les
\texttt{agent\_instructions}.

% ─────────────────────────────────────────────────────────────────────────────
\section{Séparation CalculationAgent / WriterAgent}
% ─────────────────────────────────────────────────────────────────────────────

\subsection{Deux tâches cognitives distinctes}

L'agent monolithique précédent (\texttt{WriterAgent}) mêlait deux modes
cognitifs fondamentalement différents :

\begin{table}[h]
\centering
\begin{tabularx}{\textwidth}{lXX}
\toprule
 & CalculationAgent & WriterAgent\\
\midrule
Question centrale & Les chiffres sont-ils justes ? & Le rapport est-il bon ?\\
Raisonnement & Actuariel, rigoureux, déterministe & Éditorial, narratif, qualitatif\\
Tools & \texttt{builder}, \texttt{statistical\_analysis}, \texttt{graphs}, \texttt{reasoning} & Tous (\texttt{build\_pdf}, \texttt{graphs}, \texttt{builder} en rappel)\\
Catalogue & MIDDLE / FULL / LIGHT selon phase & FULL systématiquement\\
Séquence & Quality gates bloquants & Structure narrative du rapport\\
\bottomrule
\end{tabularx}
\caption{Comparaison des deux agents}
\end{table}

\subsection{Interface de communication : le data\_store}

Le \texttt{data\_store} est le dictionnaire de résultats accumulés au fil
de l'exécution. Il constitue l'interface de communication entre les deux
agents : le \texttt{CalculationAgent} le remplit, le \texttt{WriterAgent}
le consomme. Le \texttt{WriterAgent} peut cependant rappeler des tools de
calcul s'il constate qu'une donnée manque pour compléter son rapport.

\subsection{Mécanisme de handoff}

La transition du \texttt{CalculationAgent} vers le \texttt{WriterAgent} est
signalée par l'inclusion du marqueur \texttt{<HANDOFF\_WRITER>} dans la
réponse de l'agent. Une arête conditionnelle du graphe LangGraph détecte
ce marqueur et route l'exécution vers le nœud Writer.

% ─────────────────────────────────────────────────────────────────────────────
\section{Migration vers LangGraph}
% ─────────────────────────────────────────────────────────────────────────────

\subsection{Pourquoi LangGraph maintenant ?}

Trois facteurs ont motivé le choix de migrer vers LangGraph dès cette
refactorisation plutôt qu'attendre :

\begin{enumerate}
  \item La gestion du niveau de catalogue (MIDDLE / FULL / LIGHT) crée un
    état qui doit être porté entre les nœuds. LangGraph fournit ce mécanisme
    nativement via le \texttt{StateGraph} et le \texttt{AgentState}.

  \item La séparation CalculationAgent / WriterAgent implique un mécanisme
    de routing conditionnel que LangGraph rend explicite et traçable.

  \item L'agent va devoir couvrir d'autres domaines actuariels (tarification,
    provisionnement). L'ajout d'un superviseur LangGraph pour router entre
    agents spécialisés sera trivial à partir de cette architecture, alors qu'il
    nécessiterait une réécriture depuis la boucle manuelle.
\end{enumerate}

\subsection{Limites de LangGraph}

Pour mémoire, LangGraph rend le code moins lisible pour les cas simples :

\begin{lstlisting}[language=Python, caption={Boucle manuelle (ancienne approche) — très lisible}]
while steps < MAX_STEPS:
    response = llm.chat.completions.create(messages=messages)
    if response has tool_call:
        result = call_tool(tool_name, params)
        messages.append(result)
    else:
        yield final_message; break
\end{lstlisting}

LangGraph exige la connaissance de \texttt{StateGraph}, \texttt{TypedDict},
\texttt{add\_conditional\_edges}, etc. La règle retenue est la suivante :
\emph{dès que la logique de l'agent nécessite un diagramme pour être expliquée,
LangGraph la rend plus lisible que la boucle manuelle}.

\subsection{Structure du graphe}

\begin{figure}[h]
\centering
\begin{tikzpicture}[
  node/.style={rectangle, draw, rounded corners, minimum width=3.5cm, minimum height=0.7cm, fill=backcolour},
  decision/.style={diamond, draw, minimum width=2.5cm, minimum height=0.7cm, fill=yellow!20, aspect=2.5},
  arrow/.style={->, thick, >=Latex},
  label/.style={font=\footnotesize, midway, above}
]

\node[node] (start)   at (0, 8)    {START};
\node[node] (router)  at (0, 7)    {router};
\node[node] (calc)    at (-3, 5.5) {calculation\_node};
\node[node] (writer)  at (3, 5.5)  {writer\_node};
\node[node] (toolsc)  at (-3, 3.5) {tools\_node (calc)};
\node[node] (toolsw)  at (3, 3.5)  {tools\_node (writer)};
\node[node] (end)     at (0, 1.5)  {END};

\draw[arrow] (start) -- (router);
\draw[arrow] (router) -- node[label, left]{calculation} (calc);
\draw[arrow] (router) -- node[label, right]{writer} (writer);

\draw[arrow] (calc) -- node[label, left]{tools} (toolsc);
\draw[arrow] (toolsc) -- (calc);
\draw[arrow] (calc) -- node[label, above, sloped]{to\_writer} (writer);
\draw[arrow] (calc) -- node[label, right, sloped]{end} (end);

\draw[arrow] (writer) -- node[label, right]{tools} (toolsw);
\draw[arrow] (toolsw) -- (writer);
\draw[arrow] (writer) -- node[label, above, sloped]{end} (end);

\end{tikzpicture}
\caption{Graphe LangGraph de l'agent actuariel}
\end{figure}

\subsection{L'AgentState}

L'état partagé entre tous les nœuds contient :

\begin{lstlisting}[language=Python, caption={AgentState — état partagé du graphe}]
class AgentState(TypedDict):
    messages:          Annotated[list[AnyMessage], add_messages]
    df_json:           str | None       # DataFrame sérialisé
    data_store:        dict             # résultats accumulés
    context_docs:      list[dict]       # documents uploadés
    plan_established:  bool             # True après le 1er tool call
    active_agent:      str              # "calculation" | "writer"
    events:            list[dict]       # events canvas (rétrocompatibilité)
    step_by_step:      bool
    pending_tool_call: dict | None
\end{lstlisting}

\subsection{Rétrocompatibilité avec l'interface canvas}

L'interface Dash (\texttt{canvas\_app.py}) attendait des événements sous forme
de dictionnaires Python générés par un générateur. Pour préserver cette interface
sans réécrire le canvas, un adaptateur \texttt{stream\_agent()} a été créé. Il
expose exactement la même signature que l'ancien \texttt{writer\_agent.run\_agent\_loop()}
et traduit les mises à jour d'état LangGraph en événements canvas identiques :

\begin{lstlisting}[language=Python, caption={Interface identique à l'ancienne API}]
# Avant
for event in WriterAgent().run_agent_loop(history, df=df, ...):
    process(event)

# Après (même interface, même events)
for event in stream_agent(history, df=df, ...):
    process(event)
\end{lstlisting}

% ─────────────────────────────────────────────────────────────────────────────
\section{Roadmap : superviseur multi-domaines}
% ─────────────────────────────────────────────────────────────────────────────

\subsection{Architecture cible future}

La refactorisation actuelle pose les fondations d'une architecture
multi-domaines. Lorsqu'un second domaine actuariel sera ajouté (tarification,
provisionnement), il suffira d'ajouter un \texttt{SupervisorAgent} LangGraph
qui route entre les agents spécialisés selon l'intent de la demande :

\begin{figure}[h]
\centering
\begin{tikzpicture}[
  node/.style={rectangle, draw, rounded corners, minimum width=3.5cm, minimum height=0.7cm, fill=backcolour},
  arrow/.style={->, thick, >=Latex}
]

\node[node] (user)    at (0, 6)    {Demande utilisateur};
\node[node] (super)   at (0, 4.5)  {SupervisorAgent};
\node[node] (mort)    at (-4, 3)   {MortalityAgent};
\node[node] (tarif)   at (-1.3, 3) {PricingAgent};
\node[node] (prov)    at (1.3, 3)  {ProvAgent};
\node[node] (writer)  at (4, 3)    {WriterAgent};

\draw[arrow] (user) -- (super);
\draw[arrow] (super) -- (mort);
\draw[arrow] (super) -- (tarif);
\draw[arrow] (super) -- (prov);
\draw[arrow] (super) -- (writer);
\draw[arrow] (mort) -- (writer);
\draw[arrow] (tarif) -- (writer);
\draw[arrow] (prov) -- (writer);

\end{tikzpicture}
\caption{Architecture cible multi-domaines avec superviseur LangGraph}
\end{figure}

\subsection{Ce qui sera trivial à migrer}

\begin{itemize}
  \item L'actuel \texttt{CalculationAgent} devient le \texttt{MortalityAgent}
    sans modification.
  \item Le \texttt{WriterAgent} reste transverse à tous les domaines.
  \item Le \texttt{SupervisorAgent} reçoit le catalogue MIDDLE de chaque domaine
    pour router correctement — quelques centaines de tokens supplémentaires.
  \item Le \texttt{data\_store} devient l'interface contractuelle entre domaines :
    un agent de tarification peut consommer l'exposition produite par
    le MortalityAgent.
\end{itemize}

% ─────────────────────────────────────────────────────────────────────────────
\section{Bilan de la refactorisation}
% ─────────────────────────────────────────────────────────────────────────────

\subsection{Fichiers créés et modifiés}

\begin{table}[h]
\centering
\begin{tabularx}{\textwidth}{lXl}
\toprule
Fichier & Rôle & Action\\
\midrule
\texttt{report\_agent/agents/state.py} & \texttt{AgentState} TypedDict & Créé\\
\texttt{report\_agent/agents/tools\_node.py} & Exécution tools + step-by-step & Créé\\
\texttt{report\_agent/agents/calculation.py} & Nœud CalculationAgent & Créé\\
\texttt{report\_agent/agents/writer.py} & Nœud WriterAgent & Créé\\
\texttt{report\_agent/agents/graph.py} & Graphe LangGraph + \texttt{stream\_agent()} & Créé\\
\texttt{report\_agent/tools/catalogue.py} & 3 méthodes MIDDLE/FULL/LIGHT & Modifié\\
\texttt{loader.py} & Paramètre \texttt{level} & Modifié\\
\texttt{canvas\_app.py} & Worker thread → \texttt{stream\_agent()} & Modifié\\
\texttt{run\_pipeline.py} & → \texttt{stream\_agent()} & Modifié\\
\texttt{requirements.txt} & \texttt{langgraph}, \texttt{langchain-openai} & Modifié\\
\bottomrule
\end{tabularx}
\caption{Bilan des fichiers impactés par la refactorisation}
\end{table}

\subsection{Résultats du pipeline de validation}

Le pipeline complet a été exécuté sur le portefeuille de référence
(530~345 lignes) après la refactorisation, avec les résultats suivants :

\begin{itemize}
  \item \textbf{40 étapes exécutées} sans crash ni erreur 429
  \item \textbf{Lissage adaptatif} : Whittaker escaladé de $\lambda=1000$ à
    $\lambda=50000$, puis basculement automatique sur Gompertz
    (\texttt{n\_non\_monotone = 0}, $R^2 = 0.93$)
  \item \textbf{Validation} : 46 âges avec intervalles de confiance Poisson
    (largeur moyenne 0.023)
  \item \textbf{Benchmarking} : facteur d'abattement global 0.496 vs TH0002
  \item \textbf{PDF généré} avec succès
  \item \textbf{Prompt système} : 4~553 tokens sur les tours d'exécution
    (vs 17~000 auparavant, soit $-73\%$)
\end{itemize}

\end{document}
